%----------------------------------------------------------------------------
\chapter{Introduction}
\label{chap:introduction}
%----------------------------------------------------------------------------

\section{Motivation}
\label{sec:introduction-motivation}

Autonomous surface vehicles (ASVs) are increasingly used for tasks such as
environmental monitoring, hydrographic surveying and maritime security, where
they can operate for long periods at lower cost and risk than crewed
vessels~\cite{TODO-asv-overview}. A fundamental requirement of every such
system is the ability to determine its own position and to build a model of
its surroundings.

This thesis is motivated by the plans of the BME Solar Boat Team to enter new
competition categories. The team's latest vessel, the trimaran
\textit{Crystal}, was designed for solar-class racing; however, the number of
solar-class competitions has been declining in recent seasons. The team is
therefore turning towards autonomous boating, with the longer-term goal of
entering the autonomous class of Monaco Energy Boat Challenge in which full-scale
boats race without a crew.

RoboBoat, an international student competition organized by
RoboNation~\cite{TODO-roboboat-handbook}, offers a suitable entry point. Its
participants build small ASVs that must complete a series of autonomous tasks
on a course marked by buoys. Most of these tasks require the vehicle to detect
and classify buoys and to navigate relative to them, which in turn requires a
consistent estimate of both the vehicle's position and the buoy locations.
Because RoboBoat vehicles are small and inexpensive, the competition also
provides a low-risk platform for developing autonomy systems before deploying
them on a full-scale boat.

The focus of this thesis is the localization and mapping subsystem of such a
vehicle: estimating the pose of the ASV while simultaneously building a map of
the classified buoys around it.


\section{Problem Statement}
\label{sec:introduction-problem-statement}
% TODO: Define the core problem addressed by this thesis
An ASV competing in RoboBoat operates on a course whose buoy positions are not
known in advance. GNSS alone provides the vehicle's position with meter-level
error, carries no information about the surrounding buoys, and may degrade
near shore structures. Conversely, buoy detections from onboard sensors are
only useful if the vehicle's own pose is known accurately. Perception on water
adds further difficulty: small buoys are sparsely sampled by LiDAR at longer
ranges, while reflections and glare corrupt both LiDAR and camera data. The
RoboBoat designs reviewed in this work typically rely on GNSS combined with
task-level buoy tracking, and do not report the accuracy or consistency of
their estimates~\cite{TODO}.

The problem addressed in this thesis is therefore how the pose of an ASV and
the map of classified buoys can be estimated jointly and consistently by
fusing GNSS, inertial, LiDAR and camera measurements, and how this estimate
behaves under varying sensor noise, wave conditions and GNSS outages.


\section{Objectives and Scope}
\label{sec:introduction-objectives}
% TODO: Simulation only; control and path planning out of scope



\section{Structure of the Thesis}
\label{sec:introduction-structure}
% TODO: Brief overview of each chapter


